\begin{solapartat}
\punts{0,6} $\phi = \vec{B}\cdot\vec{S} = B\cdot S\cdot\cos(0) = 0,1\cdot 0,02\cdot 0,016\cdot 1 = 3,20\times 10^{-5}\ \mathrm{Wb}$

\punts{0,65} $\phi = \vec{B}\cdot\vec{S} = B\cdot S\cdot\cos(\theta) = 0,1\cdot 0,02\cdot 0,016\cdot\cos(\theta) = 3,20\times 10^{-5}\cos(\theta)\ \mathrm{Wb}$
\end{solapartat}

\begin{solapartat}
\punts{0,2} L'angle en funció del temps en segons: $\theta = \omega t = 60\ \dfrac{\text{voltes}}{\text{s}}\cdot\dfrac{2\pi}{1\ \text{volta}}\cdot t = 120\pi\, t\ \mathrm{rad}$

\punts{0,3} Tenint en compte les condicions inicials, el flux total a través de la bobina és:
\[ \phi = N\cdot\vec{B}\cdot\vec{S} = 100\cdot 3,20\times 10^{-5}\cos(120\pi\, t) = 3,20\times 10^{-3}\cos(120\pi\, t)\ \mathrm{Wb} \]

\punts{0,5} I a partir de la llei de Faraday, la fem induïda és:
\[ \varepsilon = -\frac{d\phi}{dt} = -\frac{d}{dt}\left[3,20\times 10^{-3}\cos(120\pi\, t)\right] = 120\pi\cdot 3,20\times 10^{-3}\sin(120\pi t) \]
I la força electromotriu s'expressa com:
\[ \varepsilon = 1,21\sin(120\pi t)\ \mathrm{V}\ (t\ \text{en s}) \]

\punts{0,25} Atès que $\varepsilon = \varepsilon_\text{màx}\cdot\sin(\omega t)$, llavors la força electromotriu màxima és:
\[ \varepsilon_\text{màx} = 1,21\ \mathrm{V} \]
\end{solapartat}
